\apphead{Appendix 3}{Hecke Polynomials, coherent cohomology, and U}

In this \uline{final} appendix, we explain the relation between Hecke
polynomials mod p, coherent cohomology, and the endomorphism $U$ of
$S(R,r,n,0)\tens K$ (notations as in \eqref{ka:3.11}).

\sechead{A3.1}{The Fredholm determinant of U}

\kaitem{A3.1.0} Let us begin by computing the \uline{trace} of $U^n$, using
the Dwork-Monsky fixed point formula. For simplicity, we take $R$ to have
residue field $\F_p$. Let $R_m$ be its unramified extension of degree $m$,
and $K_m$ the fraction field of $R_m$. The endomorphism $\varphi$ acts on the
points of $\Mbar_n(\F_p,1)$ with values in the algebraic closure
$\overline{\F}_p$ of $\F_p$ as the relative Frobenius\index{Frobenius}. For each integer
$m\geq 1$ we denote by $T_m^0$ the set of $\overline{\F}_p$-valued points of
$\Mbar_n(\F_p,1)$ which are \uline{fixed} by the m'th iterate $\varphi^m$,
i.e., $T_m^0$ is the set of $\F_{p^m}$-valued points of $\Mbar_n(\F_p,1)$.
It is known (cf.\cite{monsky71}) that each element of $T_m^0$ lifts to a unique
$R_m$-valued point of the \uline{formal scheme} $\Mbar_n(R,1)$ which is fixed
by $\varphi$. We denote by $T_m$ the set of such $\varphi$-fixed $R_m$-valued
points of $\Mbar_n(R,1)$ (so $T_m\xrightarrow{\sim}T_m^0$ by reduction mod
$\cP$). The tangent space to $\Mbar_n(R,1)$ at a point $t\in T_m$ is a free
$R_m$-module of rank one, on which $\varphi^m$ acts as an $R_m$-linear
endomorphism. We denote by $d\varphi^m(t)\in R_m$ its "matrix". The
Dwork-Monsky trace formula \cite{monsky71} is as follows:

\begin{kadisp}{A3.1.1}
\operatorname{trace}(U^m) = \sum_{t\in T_m}\frac{-1}{p^m}\frac{d\varphi^m(t)}{1-d\varphi^m(t)} .
\end{kadisp}

It remains to determine the "local terms" in this formula. We begin with the
cusps, i.e., the points $t\in T_m$ whose image $t_0\in T_m^0$ is a cusp of
$\Mbar_n(\F_q,1)$. Then, as we have seen, the overlying point $t\in T_m$ is
\uline{itself} a cusp of $\Mbar_n(R_m,1)$, the completion of its local ring is
$R_m[[q]]$, and the action of $\varphi^m$ is given by
$q\longrightarrow q^{p^m}$, whose \uline{linear} term is \uline{zero}. Hence
$d\varphi^m(t)=0$ at the cusps.

Now suppose $t\in T_m$ is not a cusp. Then the corresponding elliptic curve\index{elliptic curve}
$E_t$ is the so-called \uline{canonical lifting} of its reduction $E_{t_0}$
(because the m'th iterate of the Frobenius endomorphism of $E_{t_0}$ lifts to
an endomorphism of $E_t$, namely m-fold division by the canonical subgroup\index{canonical subgroup} -
(cf.Messing \cite{messing72}). In this case it is known that the completion of the local
ring at $t$ is isomorphic to $R_m[[X]]$, where $1+X$ is the Serre-Tate
parameter (cf. ft note, p.186). Let $\alpha\in\Z_p^\times=\alpha(m,t_0)$ be
the "matrix" of the action of the automorphism "$p^m$-th power" acting on the
Tate module $T_p(E_{t_0}(\overline{\F}_p))$ of the reduced curve. Then
(cf.Messing \cite{messing72}), the action of $\varphi^m$ on $R_m[[X]]$ is the one sending
$1+X\longrightarrow(1+X)^{p^m/\alpha^2}$, hence
$d\varphi^m(t)=p^m/\alpha(m,t_0)^2$. Combining all this, we find the formula

\begin{kadisp}{A3.1.2}
\operatorname{trace}(U^m) = \sum_{\substack{t\in T_m\\ t\text{ not a cusp}}}\frac{1}{p^m-\alpha(m,t_0)^2} .
\end{kadisp}

Denoting by $T_m^{oo}$ the set of $\F_{p^m}$-valued points of $M_n(\F_p,1)$,
i.e. the set of \uline{ordinary} elliptic curves over $\F_{p^m}$ with level n
structure, we have

\begin{kadisp}{A3.1.3}
\operatorname{trace}(U^m) = \sum_{t_0\in T_m^{oo}}\frac{1}{p^m-\alpha(m,t_0)^2} .
\end{kadisp}

The next step is to assemble this data into an expression for the Fredholm
determinant $\det(1-tU)$ as a product of L-series on $M_n(\F_p,1)$. For any
\uline{closed point} $x$ of $M_n(\F_p,1)$ (i.e., an \uline{orbit} of
$\Gal(\overline{\F}_p/\F_p)$ acting on the $\overline{\F}_p$-valued points of
$M_n(\F_p,1)$), we define $\alpha(x)=\alpha(\deg(x),\widetilde{x})$, where
$\widetilde{x}$ is any $\F_{p^{\deg(x)}}$-valued point of $M_n(\F_p,1)$ lying
over $x$. For each integer $r$, the L-series $L(M_n(\F_p,1);\alpha^r;t)$ is
the element of $\Z_p[[t]]$ given by the infinite product over all closed
points $x$ of $M_n(\F_p,1)$

\begin{kadisp}{A3.1.4}
\prod_x (1-\alpha^r(x)\cdot t^{\deg(x)})^{-1} .
\end{kadisp}

An elementary calculation now yields the following identity.

\begin{kathm}{Identity}{A3.1.5}
\[
\det(1-tU) = \prod_{r\geq 0} L(M_n(\F_p,1),\alpha^{-2(r+1)},p^r t)\qquad\text{(which is the key point of \cite{dwork71}).}
\]
It shows \uline{independently} of \eqref{ka:3.11.7} that $\det(1-tU)$ lies in
$\Z_p[[t]]$, and gives as a corollary the following congruence formula.
\end{kathm}

\begin{kathm}{Corollary}{A3.1.6.}
$\det(1-tU)\equiv L(M_n(\F_p,1),\alpha^{p-3},t)$ modulo $p\cdot\Z_p[[t]]$ .
\end{kathm}

\begin{kaproof}
the term with $r=0$ remains modulo $p$, and modulo $p$ the characters
$\alpha^{-2}$ and $\alpha^{p-3}$ are equal, hence give L-series which coincide
mod $p$.
\end{kaproof}

But the character $\alpha_0=\alpha$ mod $p$ is the one associated to the
locally constant rank-one $\F_p$-etale sheaf $R^1f_*\F_p$, and the L-series
$L(M_n(\F_p,1),\alpha_0^{p-3},t)$ is just the L-series
$L(M_n(\F_p,1),(R^1f_*\F_p)^{\otimes p-3},t)$ associated to
$(R^1f_*\F_p)^{\otimes p-3}$ in \eqref{ka:4.1.1}.

[NB the apparent inversion is due to the fact that $\alpha$ describes the
action of the arithmetic Frobenius on the etale quotient of Ker p, and hence
by duality it is the action of the geometric Frobenius on its dual
$R^1f_*\F_p$.]

Furthermore, the sheaf $R^1f_*\F_p$ \uline{extends} to a locally constant
rank-one $\F_p$-etale sheaf on $\Mbar_n(\F_p,1)$, and the value of the
\uline{extended} character (still denoted $\alpha_0$) is 1 at each cusp
(cf.\eqref{ka:4.2.1}). Thus we have

\begin{kadisp}{A3.1.7}
L(M_n(\F_p,1),\alpha_0^{p-3},t) = [\prod_{\substack{x\text{ closed}\\ \text{point among}\\\text{the cusps}}}(1-t^{\deg x})]\cdot L(\Mbar_n(\F_p,1),(R^1f_*\F_p)^{\otimes p-3},t) .
\end{kadisp}

\sechead{A3.2}{Relation to mod p \'etale cohomology and to coherent cohomology}

Let $H^1_{comp}$ denote the etale cohomology groups with compact
supports $H^1_{comp}(\Mbar_n(\overline{\F}_p,1),(R^1f_*\F_p)^{\otimes p-3})$,
which are $\Gal(\overline{\F}_p/\F_p)$-modules over $\F_p$. Only $H^1_{comp}$
is $\neq 0$. Let $F_{geom}\in\Gal(\F_p/\F_p)$ denote the \uline{inverse} of
the automorphism $x\longrightarrow x^p$. According to (\cite{deligne-sga4}), we have the
formula

\begin{kadisp}{A3.2.1}
L(\Mbar_n(\F_p,1),(R^1f_*\F_p)^{\otimes p-3},t) = \det(1-tF_{geom}|H^1_{comp}) .
\end{kadisp}

By \eqref{ka:4.2.2}, the invertible sheaf with p-linear "automorphism" corresponding to
$(R^1f_*\F_p)^{\otimes p-3}$ is $(\w^{3-p},\varphi)$ over $\Mbar_n(\F_p,1)$.
But the pair $(\w^{\otimes 3-p},\varphi)$ extends to an invertible sheaf with
p-linear endomorphism on \uline{all} of $\Mbar_n\otimes\F_p$, namely to the
invertible sheaf $\w^{\otimes 3-p}$ on $\Mbar_n\otimes\F_p$, with p-linear
endomorphism

\[
\overline{\varphi}\colon \w^{\otimes 3-p}\longrightarrow \w^{\otimes 3-p}
\]

given by

\[
\overline{\varphi}\colon g\longrightarrow A^{p-3}\cdot g^p
\]

where $A=E_{p-1}$ mod $p$ denotes the \uline{Hasse invariant\index{Hasse invariant}}
$\in\Gamma(\Mbar_n\otimes\F_p,\w^{\otimes p-1})$, (compare q-expansions!)

Because this \uline{extended endomorphism} \uline{vanishes} at the fibres
\uline{outside} $\Mbar_n(\F_p,1)$, we have an isomorphism

\begin{kadisp}{A3.2.2}
& H^1_{comp} \cong \text{the \uline{fixed} points of }\overline{\varphi}\text{ acting p-linearly on} \\
& H^1(\Mbar_n\otimes\F_p,\w^{\otimes 3-p})
\end{kadisp}

under which the action of the \uline{arithmetic} Frobenius\index{Frobenius} on $H^1_{comp}$ is
its obvious action on the fixed points of $\overline{\varphi}$. It follows
formally that we have the identity

\begin{kadisp}{A3.2.3}
\det(1-tF_{geom}|H^1_{comp}) = \det(1-t\overline{\varphi}|H^1(\Mbar_n\otimes\F_p,\w^{\otimes 3-p})) .
\end{kadisp}

Putting this all together, we have the following congruence relation modulo
$p\Z_p[[t]]$.

\begin{kadisp}{A3.2.4}
\det(1-tU)\equiv [\prod_{\substack{x\text{ closed}\\ \text{point lying}\\\text{among the cusps}}}(1-t^{\deg x})]\cdot\det(1-t\overline{\varphi}|H^1(\Mbar_n\otimes\F_p,\w^{\otimes 3-p})) .
\end{kadisp}

\sechead{A3.3}{Relation to the Cartier operator}

We now wish to calculate the determinant of $\overline{\varphi}$
on $H^1(\Mbar_n\otimes\F_p,\w^{\otimes 3-p})$ by using Serre duality and the
\uline{Cartier operator\index{Cartier operator}}. For this it is convenient to abstract the situation
slightly in the following lemma - in which $X$ is $\Mbar_n\otimes\F_p$,
$\cL$ is $\w^{\otimes p-3}$, and $B$ is $A^{p-3}$.

\begin{kathm}{Lemma}{A3.3.1.}
Let $X$ be a projective smooth curve over $\F_p$, $\cL$ an invertible sheaf,
and $B$ a section of $\cL^{\otimes p-1}$. The composition
\end{kathm}

\begin{kadisp}{A3.3.2}
\cL\otimes\Omega^1_X\xrightarrow{B}\cL^{\otimes p}\otimes\Omega^1_X\xrightarrow{C}\cL\otimes\Omega^1_X
\end{kadisp}

(where $C$ is the Cartier operation, defined locally by
$C(\ell^p\otimes\omega)=\ell\otimes C(\omega)$) induces an endomorphism of
$H^0(X,\cL\otimes\Omega^1)$ which is \uline{dual} to the endomorphism of
$H^1(X,\cL^{-1})$ induced by the endomorphism
$\check{\ell}\longrightarrow B(\check{\ell})^p$ of $\cL^{-1}$.

\begin{kaproof}
We begin by remarking that although $X\xrightarrow{C}\PP^n$ need not be
geometrically connected, Serre duality on $\PP^n$ gives a perfect pairing
between $H^i(X,\cF)$ and $\operatorname{Ext}^{1-i}_{\cO_X}(\cF,\Omega^1_X)$
with values in $H^n(\PP^n,\Omega^m_{\PP^n})\cong\F_p$ for any coherent $\cF$
on X, which just as in the usual case may be computed via repartitions and
residues. The desired duality now follows from the fact that if $x\in X$ is a
closed point, and $\check{\ell}$ and $\xi$ are meromorphic sections of
$\cL^{-1}$ and $\cL\otimes\Omega^1$, then
$\operatorname{residue}_x(B\cdot(\check{\ell})^p\cdot\xi) =
(\operatorname{residue}_x(\check{\ell}\cdot C(B\xi)))^p$, (the usual Cartier
formula applied to the one-form $B(\ell)^p\xi)$.
\end{kaproof}

\begin{kathm}[A3.3.3-lem]{Lemma}{A3.3.3.}
Take $X=\Mbar_n\otimes\F_p$, $\cL=\w^{\otimes 2k}$, $k\geq 0$, and $B=A^{2k}$
in the previous lemma. Under the isomorphism
\end{kathm}

\[
\begin{gathered}
H^0(\Mbar_n\otimes\F_p,\w^{\otimes 2k}\otimes\Omega^1)\cong H^0(\Mbar_n\otimes\F_p,\w^{\otimes 2k+2}\otimes I(\text{cusps}))\\
\cong \text{the space of cusp forms of level }n\text{ and weight }2k+2\text{ over }\F_p,\\
\text{the endomorphism }\xi\longrightarrow C(A^{2k}\xi)\text{ is the Hecke operator\index{Hecke operator} }T_p .
\end{gathered}
\]

\begin{kaproof}
It suffices to check the q-expansions. But in terms of q-expansions and the
isomorphism $\Omega^1_X(\log\text{ cusps})\cong\w^{\otimes 2}$, if $\xi$ in
q-expansion is $f(q)\left(\frac{\dd q}{q}\right)^{k+1}$ then $C(A^k\xi)$ in
q-expansion is
$C(f(q)\cdot\left(\frac{\dd q}{q}\right)^{\otimes(2kp+2)/2}) =
C(f(q)\frac{\dd q}{q})\cdot\left(\frac{\dd q}{q}\right)^{\otimes k}$. But if
$f(q)=\sum a_nq^n$, $C(f(q)\frac{\dd q}{q}) =
\sum (a_{np})^{1/p}\cdot q^n\frac{\dd q}{q}$. Comparing this with the
\uline{explicit} formula \eqref{ka:1.11.1.2} for $T_p$ gives the desired result
(because $p^{2k-1}\equiv 0$ (p)!).
\end{kaproof}

Putting this all together, we obtain the congruence relation mod $p\Z_p[[t]]$:

\begin{equation}\tag{A3.3.3}\label{ka:A3.3.3}
\left\{
\begin{aligned}
\det(1-tU) &\equiv \Big[\prod_{\substack{x\text{ closed}\\ \text{point lying}\\\text{among the cusps}}}(1-t^{\deg x})\Big]\\
&\qquad\cdot\det\Big(1-tT_p\Big|\substack{\text{cusp forms of weight }p-1\\ \text{and level }n}\Big) \\
\det\Big(1-tT_p\Big|\substack{\text{cusp forms of weight }p-1\\ \text{and level }n}\Big) &\equiv \det\Big(1-t\cdot C\cdot A^{p-3}\Big|\\
&\qquad\qquad H^0(\Mbar_n\otimes\F_p,\w^{\otimes p-3}\otimes\Omega^1)\Big)
\end{aligned}
\right.
\end{equation}

This formula is the \uline{starting point} for recent work of Adolphson \cite{adolphson73}.

\begin{kathmstar}{FOOTNOTE :}
the first new sentence on page 182 is \uline{incorrect}, though the tangent
calculation we deduce from it is correct. The difficulty is that the
Serre-Tate parameter is \uline{not} "rational" over $R_m$, but only over
$R_\infty$, the completion of the maximal unramified extension of R. However,
\uline{if} we view t as defining, by extension of scalars, a rational point of
$\Mbar_n(R_\infty,1)$, \uline{then} the completion of its local ring
\uline{is} indeed isomorphic to $R_\infty[[X]]$, where $1+X$ is the Serre-Tate
parameter (cf. Messing \cite{messing72}). Further, the $R_\infty$-linear endomorphism of
$R_\infty[[X]]$ \uline{deduced} from $\varphi^m$ by extension of scalars
\uline{is} given by $1+X\longleftrightarrow (1+X)^{p^m/\alpha^2}$, in the
notation of page 182, and the formula \eqref{ka:A3.1.2} remains true.
\end{kathmstar}
